\documentclass[12pt]{article}
\usepackage[utf8]{inputenc}
\usepackage[T1]{fontenc}
\usepackage{hyperref}
\title{Tok Pisin Language Resources for K-12 Curriculum Development}
\author{Compiled by Claude Sonnet 5 (claude-sonnet-5) for Andr\'as Kornai}
\date{}
\begin{document}
\maketitle
\section{Introduction}
Tok Pisin (ISO 639-3 tpi; Glottocode tokp1240) is an English-based creole. It is an official national language of Papua New Guinea. About 4 million people speak it, mostly as a second language. The searches were limited. They found university course pages, a newspaper, general references and policy studies. They found no school textbooks.
\section{Primary Educational Resources}
No Tok Pisin school textbooks were found in search results. Papua New Guinea moved to community-language elementary education in the 1990s. A 2016 presentation describes an elementary model that uses Tok Pisin or a vernacular in Prep to Grade 2 and then moves gradually to English. One study from the Western Highlands reports a long-standing directive to minimise Tok Pisin in school.
\section{Linguistic and Reference Materials}
\begin{itemize}
\item Wikipedia article on Tok Pisin: \url{https://en.wikipedia.org/wiki/Tok_Pisin}
\item ANU beginner Tok Pisin course page. It lists the Oxford \emph{Papua New Guinea Tok Pisin English Dictionary} as required reading. Some editions also list Dutton's \emph{A New Course in Tok Pisin} and the \emph{Jacaranda Dictionary and Grammar of Melanesian Pidgin}: \url{https://programsandcourses.anu.edu.au/2024/course/TOKP1002/First%20Semester/3167}
\item APiCS Online contribution on Tok Pisin: \url{https://apics-online.info/contributions/22}
\end{itemize}
\section{Religious and Community Resources}
The search results mention a Tok Pisin Bible (\emph{Buk Baibel long Tok Pisin}) only as a link label. I did not see its own URL.
\section{Digital and Audio Resources}
\begin{itemize}
\item Wantok Niuspepa, a weekly newspaper that is mainly in Tok Pisin but also carries English: \url{https://en.wikipedia.org/wiki/Wantok_Niuspepa}
\end{itemize}
\section{Curriculum Development Context}
\begin{itemize}
\item UNESCO, ``Vernacular education in Papua New Guinea'': \url{https://unesdoc.unesco.org/ark:/48223/pf0000146660}
\item Devpolicy 2016 presentation on bilingual education in Papua New Guinea: \url{https://devpolicy.org/Events/2016/PNG-Update/1b_Devette-Chee.pdf}
\item ANU study of language ecology, policy and pedagogy in a Highland community: \url{https://openresearch-repository.anu.edu.au/items/2f70cde3-c473-4dda-9982-6e6220a8a6f7}
\end{itemize}
The current official language-of-instruction policy should be checked against Department of Education documents. The search results did not give a current statement.
\end{document}